% problem_id: S046-ct:yoneda
% source_id: S046 (HoTT Book)
% source_file: categories.tex
% statement_locator: categories.tex lines 941--948
% proof_locator: categories.tex lines 949--971
% env: thm
% pairing: tight (verified)
% status: complete (statement + author proof verbatim + reason + locators)
%
\begin{thm}[The Yoneda lemma]\label{ct:yoneda}
  \indexdef{Yoneda!lemma}
  For any precategory $A$, any $a:A$, and any functor $F:\uset^{A\op}$, we have an isomorphism
  \begin{equation}\label{eq:yoneda}
    \hom_{\uset^{A\op}}(\y a, F) \cong Fa.
  \end{equation}
  Moreover, this is natural in both $a$ and $F$.
\end{thm}

\begin{proof}
  Given a natural transformation $\alpha:\y a \to F$, we can consider the component $\alpha_a : \y a(a) \to F a$.
  Since $\y a(a)\jdeq \hom_A(a,a)$, we have $1_a : \y a(a)$, so that $\alpha_a(1_a) : F a$.
  This gives a function $(\alpha \mapsto \alpha_a(1_a))$ from left to right in~\eqref{eq:yoneda}.

  In the other direction, given $x:F a$, we define $\alpha:\y a \to F$ by
  \[\alpha_{a'}(f) \defeq F_{a,a'}(f)(x). \]
  Naturality is easy to check, so this gives a function from right to left in~\eqref{eq:yoneda}.

  To show that these are inverses, first suppose given $x:F a$.
  Then with $\alpha$ defined as above, we have $\alpha_a(1_a) = F_{a,a}(1_a)(x) = 1_{F a}(x) = x$.
  On the other hand, if we suppose given $\alpha:\y a \to F$ and define $x$ as above, then for any $f:\hom_A(a',a)$ we have
  \begin{align*}
    \alpha_{a'}(f)
    &= \alpha_{a'} (\y a_{a,a'}(f)(1_a))\\
    &= (\alpha_{a'}\circ \y a_{a,a'}(f))(1_a)\\
    &= (F_{a,a'}(f)\circ \alpha_a)(1_a)\\
    &= F_{a,a'}(f)(\alpha_a(1_a))\\
    &= F_{a,a'}(f)(x).
  \end{align*}
  Thus, both composites are equal to identities.
  We leave the proof of naturality to the reader.
\end{proof}
